\documentclass[border=10pt]{standalone}
\usepackage{tikz}
\usetikzlibrary{arrows.meta}

\begin{document}
\begin{tikzpicture}[
    >=Stealth,
    node distance = 3cm,
    every node/.style = {
        circle,
        draw=teal,
        fill=blue!10,
        thick,
        text=blue,
        align=center,
        font=\sffamily,
        minimum size=4cm
    },
    grayline/.style = {gray, thin},
    bluearrow/.style = {->, blue, thick}
]

% Define coordinates
\coordinate (A) at (90:3cm);
\coordinate (B) at (30:3cm);
\coordinate (C) at (-30:3cm);
\coordinate (D) at (-90:3cm);
\coordinate (E) at (-150:3cm);
\coordinate (F) at (-210:3cm);

% Draw gray lines between all pairs
\foreach \p1/\p2 in {A/B, A/C, A/D, A/E, A/F, B/C, B/D, B/E, B/F, C/D, C/E, C/F, D/E, D/F, E/F} {
    \draw[grayline] (\p1) -- (\p2);
}

% Draw nodes
\node at (A) {market assessment};
\node at (B) {objective setting};
\node at (C) {resource allocation};
\node at (D) {strategic planning};
\node at (E) {execution \& monitoring};
\node at (F) {review \& adaptation};

% Draw blue arrows in cycle
\draw[bluearrow] (A) -> (B);
\draw[bluearrow] (B) -> (C);
\draw[bluearrow] (C) -> (D);
\draw[bluearrow] (D) -> (E);
\draw[bluearrow] (E) -> (F);
\draw[bluearrow] (F) -> (A);

\end{tikzpicture}
\end{document}